\documentclass[10pt,a4paper]{amsart}
\usepackage[margin=19mm]{geometry}
\usepackage{amsmath,amssymb,mathtools}
\usepackage[scr=rsfs,cal=euler]{mathalfa}
\usepackage{xfrac}
\usepackage[unicode,hidelinks,bookmarks=false]{hyperref}
\newcommand{\ie}{i.e.\ }
\newcommand{\cf}{cf.\ }
\newcommand{\resp}{respectively\ }
\newcommand{\Subsection}{\subsection}
\providecommand{\nobreakdash}{\nobreak\hbox{-}}
\setlength{\emergencystretch}{2em}
\multlinegap0pt
\usepackage{tikz}
\usepackage{amssymb}
\usepackage{bm}
\usepackage{algorithm}\usepackage{algpseudocode}
\newtheorem{lem}{Lemma}
\newtheorem{thm}{Theorem}
\usepackage{cleveref}
\crefname{lem}{Lemma}{Lemmas}
\crefname{thm}{Theorem}{Theorems}
\newcommand{\E}{\mathbb{E}}
\newcommand{\Eerrw}{\E_{\textnormal{errw}}}
\renewcommand{\P}{\mathbb{P}}
\newcommand{\Perrw}{\P_{\textnormal{errw}}}
\newcommand{\cov}{\textnormal{cov}}
\newcommand{\diam}{\textnormal{diam}}
\newcommand{\mc}{\mathcal{C}}
\newcommand{\mr}{\mathcal{R}}
\newcommand{\poly}{\textnormal{poly}}
\title{On Statistical Estimation of Edge-Reinforced Random Walks}
\author{Qinghua (Devon) Ding, Venkat Anantharam}
\date{Source: arXiv:2503.06115v2 (CC BY 4.0)}
\begin{document}
\maketitle

\noindent\emph{Source and licence.} On Statistical Estimation of Edge-Reinforced Random Walks. Version arXiv:2503.06115v2 (CC BY 4.0), distributed by arXiv under the Creative Commons Attribution 4.0 International licence, http://creativecommons.org/licenses/by/4.0/, as recorded in the arXiv licence metadata for this exact version and shown on the abstract page https://arxiv.org/abs/2503.06115v2. Full licence notice: \texttt{licenses/arxiv-2503.06115-CC-BY-4.0.txt}.\par\smallskip

\noindent\textit{Editorial scope.} Counted unit: the article's thm thm:bound\_cover (source0.tex lines 1790--1807). The article's complete original proof of it is delivered with it (source0.tex lines 1810--2166, one \texttt{proof} environment of 357 lines). No mathematics is written, repaired or re-derived: the statement and the proof are the article's own lines, in the article's own notation, and no retained line is substituted or re-flowed.  Retained uncounted as the source-local context the proof needs: lines 1236--1240; lines 1356--1365; lines 1545--1563; lines 1622--1641; lines 1725--1729; lines 1748--1759.  Nothing was omitted from any retained span, so the package carries no omission record.  No retained line contains a citation, so the wrapper carries no bibliography and no bibliography entry.  No retained line contains a figure environment, an image-inclusion command or drawing code, so no image is delivered and the package declares no asset dependency.  Editorial formatting only, in addition: an AMS \texttt{amsart} preamble that restates the article's own declarations and macros used by the retained lines, namely the environments \texttt{lem}, \texttt{thm} on separate counters, together with the standard packages those lines need.  The retained environments print in this document as Lemma 1, Theorem 1 in order of appearance, whereas the article prints the same items with the numbering of its own full document.  The complete native source of the article as distributed in the arXiv e-print for this exact version is bundled unchanged as \texttt{sources/174435/source0.tex}.
\begin{equation}\label{eqn:psi}
\psi_{i}:=\frac12\sum_{j:j\sim i}\beta_{ij}e^{\phi_j-\phi_i}\quad (i\neq v_0),
\qquad
\psi_{v_0}:=\frac12\sum_{j:j\sim v_0}\beta_{v_0j}e^{\phi_j-\phi_{v_0}}+\gamma.
\end{equation}

\begin{lem}[Pathwise telescoping bound]\label{lem:path-telescope}
Fix a connected simple graph $G=(V,E)$, a root $v_0\in V$, and let $(\beta,\phi,\psi)$ be the fields defined as in \eqref{eqn:psi} (with the independent $\gamma\sim\Gamma(\tfrac12,1)$ at $v_0$). For each $v\in V$, fix a shortest path
\[
P(v):\quad v_0=v_0,\,v_1,\,\ldots,\,v_k=v,
\]
and write $|P(v)|:=k$. Then
\begin{equation}\label{eq:path-tele}
e^{-\phi_v}\le 2^{\,|P(v)|}\,\sqrt{\frac{\psi_v}{\psi_{v_0}}}\,\prod_{e\in P(v)}\tilde\beta_e^{-1}.
\end{equation}
\end{lem}

\begin{thm}[Uniform upper bound for $e^{-\phi}$]\label{thm:uniform-upper-clean}
Assume the setup preceding Lemma~\ref{lem:path-telescope} and that $0<\underline{a}\le a_e\le \overline{a}<\infty$ for all $e\in E$. Let $d_{\max}:=\max_{v\in V}\deg(v)$. Then for any $\alpha>0$, with probability at least $1-4n^{-\alpha}$,
\begin{equation}
\begin{aligned}
\max_{v\in V} e^{-\phi_v}
\le 
2^{\diam}\,n^{\alpha} & \sqrt{2(1+\alpha)\log n+2(\log 2)\,\overline{a}\,d_{\max}}\, \\
& \exp\Big(
\diam\cdot\Lambda+ \sqrt{2(1+\alpha)\,\Psi_1(\underline{a}/2)\,\diam\cdot\log n}+ \frac{4}{\underline{a}}(1+\alpha)\log n
\Big), \label{eq:uniform-upper-clean}
\end{aligned}
\end{equation}
where
\[
\Lambda:= -\Psi(\underline{a})+\log\Big(\frac{\overline{a}\,d_{\max}+1}{2}\Big).
\]

Note that the upper bound can be simplified as  $\max_{v\in V} e^{-\phi_v} \le \poly(n) \, d_{\max}^{\,O(\diam)}$, assuming $\alpha, a, b$ to be constant.
\end{thm}

\begin{lem}[Exponential moment]
\label{lem:unit}
Fix $G=(V,E)$, edge parameters $A=(a_e)_{e\in E}$, and a fixed edge field $\beta$.
For any $v_0,j\in V$ and any real $t$ for which the integrals below are finite, one has
\[
\int_{\mathbb{R}^{V\setminus\{v_0\}}} e^{\,t\phi_j}\,p_\beta^{(v_0)}(\phi)\,d\phi_{V\setminus\{v_0\}}
\;=\;
\int_{\mathbb{R}^{V\setminus\{j\}}} e^{\,(1-t)\phi_{v_0}}\,p_\beta^{(j)}(\phi)\,d\phi_{V\setminus\{j\}}.
\]
Equivalently,
\[
\E_{p_\beta^{(v_0)}}\big[e^{\,t\phi_j}\big]
=
\E_{p_\beta^{(j)}}\big[e^{\,(1-t)\phi_{v_0}}\big].
\]
In particular, at $t=1$,
\[
\E_{p_\beta^{(v_0)}}\big[e^{\phi_j}\big]=1.
\]
\end{lem}

\begin{lem}[Conductance bound on the cover time]
\label{lem:cond_cover}
Given a connected graph $G=(V=[n], E)$ with positive edge conductances $(q_e)_{e\in E}$, we have
\[\frac{1}{2}\mc\mr\leq t_\cov\leq \mc\mr \log n.\]
\end{lem}

\begin{lem}[Comparison of resistive networks]
\label{lem:comp_res}
For two resistive networks on the same graph with respective conductances \( (q_e)_{e \in E} \) and \( (q_e')_{e \in E} \), if
\[
\frac{q_e}{q_e'} \in [\underline{q}, \overline{q}] \quad \forall e \in E,
\]
then
\[
\frac{\mathcal{R}'}{\mathcal{R}} \in [\underline{q}, \overline{q}],
\]
where \( \mathcal{R} \) and \( \mathcal{R}' \) are the respective maximum effective resistances.
\end{lem}

\begin{thm}[Upper bound on the cover time]
\label{thm:bound_cover}
Let $G=(V=[n],E)$ be a connected, weighted graph with positive edge weights
$A=(a_e)_{e\in E}$ and diameter $\diam$. Let
\[
d_{\max}:=\max_{v\in V}\deg(v).
\]
Assume that each $a_e$ satisfies $\underline a\le a_e\le \overline a$, where
$\underline a,\overline a$ are positive constants. There exists a constant \(g_1(\underline a,\overline a)>0\), depending only on
\(\underline a,\overline a\), such that for any \(\delta\in(0,1)\), the random environment $Q$ coupled with the ERRW satisfies
\[
t_\cov
\le
n^3\log n
\left(\frac{d_{\max}}{\delta}\right)^{2g_1(\underline a,\overline a)\diam}
\]
with probability at least $1-\delta$.
\end{thm}

\begin{proof}
If $d_{\max}=1$, then, since $G$ is connected, $G$ consists of a single edge, and the claim is immediate. We therefore assume $d_{\max}\ge2$. We repeatedly use
\[
n\le 1+d_{\max}+\cdots+d_{\max}^{\diam}\le 2d_{\max}^{\diam},
\qquad
|E|\le \frac{nd_{\max}}{2}\le d_{\max}^{2\diam}.
\]
The first bound follows by exploring the graph from any vertex up to distance $\diam$, and the second follows from the handshaking lemma. Write $R:=d_{\max}/\delta$. Since $d_{\max}\ge2$ and $\delta\in(0,1)$, we have $R\ge2$.

We first control the random edge field $\beta$. Since $\beta_e\sim\Gamma(a_e,1)$, for $0<\epsilon<1$,
\[
\Perrw^{v_0,A}(\beta_e\le \epsilon)
=
\int_0^\epsilon \frac{1}{\Gamma(a_e)}x^{a_e-1}e^{-x}\,dx
\le
\frac{\epsilon^{a_e}}{\Gamma(a_e+1)}
\le
2\epsilon^{\underline a},
\]
where we used $e^{-x}\le1$ and $\Gamma(x)>1/2$ for $x>0$. For the upper tail, exponential Markov's inequality with $\lambda=1/2$ gives, for every $\alpha>0$,
\[
\Perrw^{v_0,A}(\beta_e\ge\alpha)
\le
\Eerrw^{v_0,A}[e^{\beta_e/2}]e^{-\alpha/2}
=
2^{a_e}e^{-\alpha/2}
\le
e^{-\alpha/2+\overline a}.
\]

Choose
\[
\epsilon_\beta:=\left(\frac{\delta}{8|E|}\right)^{1/\underline a},
\qquad
\alpha_\beta:=2\left(\overline a+\log\frac{4|E|}{\delta}\right).
\]
Then the union bound gives
\[
\Perrw^{v_0,A}(\exists e\in E:\beta_e<\epsilon_\beta)
\le
2|E|\epsilon_\beta^{\underline a}
=
\frac{\delta}{4},
\]
and
\[
\Perrw^{v_0,A}(\exists e\in E:\beta_e>\alpha_\beta)
\le
|E|e^{-\alpha_\beta/2+\overline a}
=
\frac{\delta}{4}.
\]
Thus, with probability at least $1-\delta/2$, all edges satisfy $\epsilon_\beta\le\beta_e\le\alpha_\beta$.

Next, by \cref{lem:unit} with $t=1$, $\Eerrw^{v_0, A}[e^{\phi_i}]=1$ for every $i\in V$. Hence Markov's inequality and the union bound give, for any $x>0$,
\[
\Perrw^{v_0,A}\left(\max_{i\in V}e^{\phi_i}>x\right)
\le
\frac{n}{x}.
\]
Choosing $x_\phi:=4n/\delta$, we obtain
\[
\Perrw^{v_0,A}\left(\max_{i\in V}e^{\phi_i}>x_\phi\right)
\le
\frac{\delta}{4}.
\]

For the negative side of the $\phi$-field, choose $\alpha_0:=\log(16/\delta)/\log n$, so that $4n^{-\alpha_0}=\delta/4$. Applying \cref{thm:uniform-upper-clean} with this choice gives
\[
\Perrw^{v_0,A}\left(\max_{i\in V}e^{-\phi_i}>y_\phi\right)
\le
\frac{\delta}{4},
\]
where
\[
\begin{aligned}
y_\phi
:={}&
2^{\diam}n^{\alpha_0}
\sqrt{2(1+\alpha_0)\log n+2(\log2)\overline a d_{\max}} \\
&\times
\exp\left(
\diam\Lambda
+
\sqrt{2(1+\alpha_0)\Psi_1(\underline a/2)\diam\log n}
+
\frac{4}{\underline a}(1+\alpha_0)\log n
\right),
\end{aligned}
\]
and $\Lambda:=-\Psi(\underline a)+\log((\overline a d_{\max}+1)/2)$.

By the union bound, with probability at least $1-\delta$, all four estimates hold simultaneously:
\[
\epsilon_\beta\le\beta_e\le\alpha_\beta\text{ for all }e\in E,
\qquad
\max_i e^{\phi_i}\le x_\phi,
\qquad
\max_i e^{-\phi_i}\le y_\phi.
\]
On this event, for every edge $e=\{i,j\}$,
\[
Q_e=\beta_e e^{\phi_i+\phi_j}
\le
\alpha_\beta x_\phi^2,
\qquad
Q_e=\beta_e e^{\phi_i+\phi_j}
\ge
\epsilon_\beta y_\phi^{-2}.
\]
Thus
\begin{equation}
\label{eqn:Q-real-thresholds}
\forall e\in E,\qquad
Q_e\in[\epsilon_\beta y_\phi^{-2},\,\alpha_\beta x_\phi^2].
\end{equation}


It remains to compare the four thresholds with powers of \(R=d_{\max}/\delta\).
Recall that \(R\ge2\), \(\diam\ge1\), \(n\le2d_{\max}^{\diam}\le2R^\diam\), and
\(|E|\le d_{\max}^{2\diam}\le R^{2\diam}\). We also use repeatedly that
\(\log R\ge\log2\). Define
\[
C_\epsilon:=\frac{6}{\underline a},
\qquad
C_\alpha:=\frac{2\overline a}{\log2}+10,
\qquad
C_x:=5,
\]
and
\[
C_y
:=
8
+
\frac{\log(14+2(\log2)\overline a)}{\log2}
+
\frac{|\Psi(\underline a)|+\log(\overline a+1)}{\log2}
+
\sqrt{\frac{14\Psi_1(\underline a/2)}{\log2}}
+
\frac{28}{\underline a}.
\]
Set
\[
g_1(\underline a,\overline a)
:=
\max\{C_\alpha+2C_x,\ C_\epsilon+2C_y\}.
\]

First,
\[
\log\epsilon_\beta^{-1}
=
\frac{1}{\underline a}\log\frac{8|E|}{\delta}.
\]
Since \(|E|\le R^{2\diam}\) and \(\delta^{-1}\le R\),
\[
\log\frac{8|E|}{\delta}
\le
\log 8+(2\diam+1)\log R.
\]
Moreover, \(\log 8=3\log2\le3\diam\log R\), and
\(\log R\le\diam\log R\). Therefore
\[
\log\epsilon_\beta^{-1}
\le
\frac{1}{\underline a}\bigl(3\diam\log R+2\diam\log R+\diam\log R\bigr)
=
\frac{6}{\underline a}\diam\log R.
\]
Hence $\epsilon_\beta^{-1}\le R^{C_\epsilon\diam}$.

Next,
\[
\alpha_\beta
=
2\left(\overline a+\log\frac{4|E|}{\delta}\right)
\le
2\left(\overline a+\log4+(2\diam+1)\log R\right).
\]
Using \(\diam\ge1\) and \(\log R\ge\log2\), this gives
\[
\alpha_\beta
\le
\left(\frac{2\overline a}{\log2}+10\right)\diam\log R = C_\alpha\diam\log R = \log (R^{C_\alpha\diam}) \le R^{C_\alpha\diam}.
\]

Also,
\[
x_\phi=\frac{4n}{\delta}
\le
8R^{\diam+1},
\]
because \(n\le2R^\diam\) and \(\delta^{-1}\le R\). Taking logarithms gives
\[
\log x_\phi
\le
\log8+(\diam+1)\log R
\le
3\diam\log R+\diam\log R+\diam\log R
=
5\diam\log R.
\]
Thus
\(
x_\phi\le R^{C_x\diam}.
\)

It remains to bound \(y_\phi\). Recall that
\[
\alpha_0:=\frac{\log(16/\delta)}{\log n},
\qquad
n^{\alpha_0}=\frac{16}{\delta}.
\]
Hence \((1+\alpha_0)\log n=\log n+\log(16/\delta)\). Since \(n\le2R^\diam\),
\(\delta^{-1}\le R\), \(R\ge2\), and \(\diam\ge1\), we have
\[
\log n\le2\diam\log R,
\qquad
\log(16/\delta)\le5\diam\log R,
\]
and therefore
\[
(1+\alpha_0)\log n\le7\diam\log R.
\]

Taking logarithms in the definition of \(y_\phi\), we get
\[
\begin{aligned}
\log y_\phi
\le{}&
\diam\log2+\alpha_0\log n
+\frac12\log\left(2(1+\alpha_0)\log n+2(\log2)\overline a d_{\max}\right) \\
&\quad
+\diam\Lambda
+\sqrt{2(1+\alpha_0)\Psi_1(\underline a/2)\diam\log n}
+\frac{4}{\underline a}(1+\alpha_0)\log n ,
\end{aligned}
\]
where \(\Lambda=-\Psi(\underline a)+\log((\overline a d_{\max}+1)/2)\).

We now bound the terms on the right crudely. First,
\[
\diam\log2+\alpha_0\log n
=
\diam\log2+\log(16/\delta)
\le6\diam\log R.
\]
For the logarithmic square-root prefactor, since \((1+\alpha_0)\log n\le7\diam\log R\) and \(d_{\max}\le R\),
\[
2(1+\alpha_0)\log n+2(\log2)\overline a d_{\max}
\le
14\diam\log R+2(\log2)\overline a R.
\]
Using \(\diam\log R\le R^\diam\) and \(R\le R^\diam\), the quantity inside the logarithm is at most
\((14+2(\log2)\overline a)R^\diam\). Thus, after absorbing the constant into
\(\diam\log R\), this logarithmic term is at most
\[
\left(1+\frac{\log(14+2(\log2)\overline a)}{\log2}\right)\diam\log R.
\]

Next, since \(d_{\max}\le R\),
\[
\Lambda\le |\Psi(\underline a)|+\log(\overline a+1)+\log R,
\]
and hence
\[
\diam\Lambda
\le
\left(
1+\frac{|\Psi(\underline a)|+\log(\overline a+1)}{\log2}
\right)\diam\log R.
\]
For the Gaussian fluctuation term, using \((1+\alpha_0)\log n\le7\diam\log R\),
\[
\sqrt{2(1+\alpha_0)\Psi_1(\underline a/2)\diam\log n}
\le
\sqrt{14\Psi_1(\underline a/2)}\,\diam\sqrt{\log R}
\le
\frac{\sqrt{14\Psi_1(\underline a/2)}}{\sqrt{\log2}}\,\diam\log R.
\]
Finally,
\[
\frac{4}{\underline a}(1+\alpha_0)\log n
\le
\frac{28}{\underline a}\diam\log R.
\]
Combining these estimates gives
\[
\log y_\phi\le C_y\diam\log R,
\]
where one may take
\[
C_y
:=
8
+
\frac{\log(14+2(\log2)\overline a)}{\log2}
+
\frac{|\Psi(\underline a)|+\log(\overline a+1)}{\log2}
+
\sqrt{\frac{14\Psi_1(\underline a/2)}{\log2}}
+
\frac{28}{\underline a}.
\]
Equivalently, \(y_\phi\le R^{C_y\diam}\).

We have shown
\[
\epsilon_\beta^{-1}\le R^{C_\epsilon\diam},
\qquad
\alpha_\beta\le R^{C_\alpha\diam},
\qquad
x_\phi\le R^{C_x\diam},
\qquad
y_\phi\le R^{C_y\diam}.
\]
Combining these with \eqref{eqn:Q-real-thresholds}, we obtain
\[
Q_e\le R^{(C_\alpha+2C_x)\diam},
\qquad
Q_e\ge R^{-(C_\epsilon+2C_y)\diam}.
\]
Thus, with
\[
g_1(\underline a,\overline a)
:=
\max\{C_\alpha+2C_x,\ C_\epsilon+2C_y\},
\]
we have, with probability at least \(1-\delta\),
\[
\forall e\in E,\qquad
Q_e\in
\left[
R^{-g_1(\underline a,\overline a)\diam},
R^{g_1(\underline a,\overline a)\diam}
\right].
\]

On this event, \cref{lem:comp_res} implies that, if $\mr$ is the maximum effective resistance of the unit-conductance network and $\mr'$ is the corresponding maximum effective resistance for the conductances $Q$, then $\mr'\le \mr R^{g_1\diam}$. Also, $\mc':=\sum_{e\in E}Q_e\le |E|R^{g_1\diam}$. Using \cref{lem:cond_cover}, we conclude that
\[
t_\cov
\le
\mc'\mr'\log n
\le
|E|\mr\log n\,R^{2g_1\diam}.
\]
Finally, since $|E|\le n^2$ and $\mr\le n$\footnote{Note that the maximum resistance one could get is the resistance between two endpoints of an $n$-path.}, we obtain
\[
t_\cov
\le
n^3\log n
\left(\frac{d_{\max}}{\delta}\right)^{2g_1\diam}.
\]
This proves the theorem.
\end{proof}

\end{document}
